\documentclass[16pt, oneside]{report}
\usepackage[francais]{babel}
\usepackage{textcomp}
 \usepackage[pdftex]{graphicx}
 \usepackage{grffile}  
\font\myfont=cmr12 at 40pt
\graphicspath{{./images/}{./}}  


\title{{\myfont Projet libre}}
\author{mbourlet, gfernand, fbuoro, mchoong \#e1r6}
\date{}

\begin{document}

\maketitle

\chapter{Pr\'esentation du projet}
Le but du projet est de r\'ealiser un jeu vid\'eo natif sur Android.\\
Ce jeu seras un RPG au combat tour par tour mais avec un profond aspect strat\'egique, le joueur ne pourra pas choisir ses actions \`a chaque tour, il devra \`a la place programmer une IA (voir plus loin).\\
Le projet se d\'ecomposera\ en deux parties : l'engine et le jeu

\section{Engine}
L\textquotesingle engine seras un moteur basique pour le jeu, il devras \^etre r\'eutilisable sur diff\'erents projets. Son but sera d\textquotesingle acc\'el\'erer le d\'eveloppement futur et actuel du jeu. Il se divise en plusieurs Managers.

\subsection{Game}
La classe Game est la classe centrale de l'engine. C'est elle qui g\`ere la logique du programme avec son syst\`eme de state.\\
Au lancement la classe r\'ecup\'erera la sauvegarde la plus r\'ecente avant de lancer le jeu.

\subsection{Le DataManager}
Le DataManager est celui qui g\`ere les sauvegardes de l'utilisateur. Il offre deux possibilit\'es:
\begin{itemize}
	\item Sauvegarder en local: Le fichier est stock\'e sur le t\'el\'ephone de l'utilisateur dans les fichiers de l'application.
	\item Sauvegarder sur le cloud: Si l\textquotesingle utilisateur se connecte avec son compte Google+ et nous en donne l'autorisation, le DataManager sauvegarde sur le cloud.\\
\end{itemize}

\subsection{State}
Les states sont les diff\'erents \'etats du jeu, ce sont elles qui permettront de g\'erer la navigation dans l'application.

\section{Projet}
Le projet en lui m\^eme utiliseras l\textquotesingle engine. Il devra utilis\'er toutes les technologies offerte par Google, tel que le stockage sur le Cloud, le multi-joueur, les hauts-faits etc.

\chapter{Syst\'eme d\textquotesingle IA}
Le joueur pourras contr\^oler une \'equipe de 3 personnages.\\
Chaque personnage pourras avoir de nombreuses classes (Guerrier, Mage, Soigneur)\\
Chaque classe de chaque personnage aura son propre niveau.\\
Quand un personnage augmente de niveau dans une classe, il gagne des outils pour programmer son IA li\'e \`a cette classe, ces outils peuvent \^etre d\'ecomposer en plusieurs types:
\begin{itemize}
	\item Condition: Permet de conditionner le comportement de l\textquotesingle IA
	\item Action: Cela peut \^etre des sorts, des attaques, etc...
	\item Slot: Permet de poser une nouvelle paire de conditions/actions
\end{itemize}

\section{Classe ?}
Chaque personnage poss\`ede plusieurs classes, le joueur peux decider de changer la classe de ces personnages pendant un combat, permettant de faire varier les strat\'egies en fonction de l'adversaire (phase de rage, strat\'egie sp�ciale li\'ee a un adversaire, etc...)\\\\
En fonction de la fr\'equence d\textquotesingle utilisation d'une classe dans un combat, celle ci gagne plus ou moins de pourcentage de la somme d\textquotesingle experience donn\'e par ce combat.

\chapter{Planning de r\'ealisation}
Le planning du projet seras g\'erer de mani\`ere agile (Sprint chaque semaine + r\'eunion tout les jours de 15min pour faire un bilan rapide).\\
Cependant, nous allons r\'ealiser plusieurs it\'eration du projet. Chaque it\'eration auras pour but de rendre un produit fonctionnel, consistant en une am\'elioration du precedent:
\begin{itemize}
	\item It\'eration 1: R\'ealiser un combat entre deux IA basique (Allie (programmable) / Ennemie)
	\item It\'eration 2: R\'ealiser un syst\`eme basique de classe 
	\item It\'eration 3: R\'ealiser un syst\`eme de niveau + inventaire
	\item It\'eration 4: R\'ealiser un syst\`eme d'histoire / quete basique
	\item It\'eration 5: Ajout d'IA ennemie
	\item It\'eration 6: Am\'elioration du syst\`eme d'histoire / qu\^ete (branche)\\

\end{itemize}
\`A la fin de ces iterations, nous pourrons considerer que la partie obligatoire est finie. Nous pr�voyons la r�alisation de quelques bonus pouvant �tre r�alis�s en plus:
\begin{itemize}

	\item Ajout des achievements
	\item Beaucoup d'items
	\item Plus de classes
	\item Ladder solo
	\item Multijoueur
	\item Ladder multijoueur
	\item Animation enti\'ere du personnage
\end{itemize}

\chapter{Design}

Le jeu sera r\'ealis\'e en pixel art.
Voici un exemple de nos premi\`eres r\'ealisations:\\
\\
\\
\includegraphics[width=0.3\textwidth]{character-base}
\includegraphics[width=0.3\textwidth]{character-head}
\includegraphics[width=0.3\textwidth]{character-robe}
\includegraphics[width=0.3\textwidth]{character-wand}
\\
\\
\\
\begin{center}
\includegraphics[width=0.3\textwidth]{character}
\end{center}

\chapter{Communication}
La communication se fera via un Website/devblog et sera appui\'ee par un compte twitter qui d\'etaillera l'avancement du jeu en postant des photos, vid\'eos de gameplay etc. Le but \'etant d'attirer le plus de personnes \`a s'int�resser au projet avant m\^eme la sortie du jeu. C'est notamment pour cela qu'une phase de beta test sera disponible aux utilisateurs les plus interess\'es.\\
 La beta test sera g\'er\'ee a l'aide de la google d�velopper console. Le site proposera \'egalement l'inscription \`a une newsletter. 


\chapter{Mon\'etisation}
La mon\'etisation du jeu se fera via les achats in-app et les pubs. 
Les achats in-app ne devront pas bousculer l'\'equilibre du jeu, notamment dans l'optique d'un multijoueur le plus strat\'egique possible. Nous pr\'evoyons pour le moment l'achat de packs d'xp qui permettrait de lvl up plus vite. \\
Les pubs se voudront non-intrusive pour ne pas g\^acher l'\'experience utilisateur.  Nous utiliserons des api comme admob (ou autre) (utilis\'e par Rovio, Backflip Studio, Fingersoft... )pour int\'egrer des pubs cibl\'ees et donc augmenter nos revenus.


\chapter{Elements \`a v\'erifier en fin de projet}
Il faudra verifier que les diff\'erentes it\'erations fonctionnent correctement:\\
\begin{itemize}
	\item It\'eration 1: R\'ealiser un combat entre deux IA basique (Allie (programmable) / Ennemie)
	\item It\'eration 2: R\'ealiser un syst\`eme basique de classe 
	\item It\'eration 3: R\'ealiser un syst\`eme de niveau + inventaire
	\item It\'eration 4: R\'ealiser un syst\`eme d'histoire / quete basique
	\item It\'eration 5: Ajout d'IA enemie
	\item It\'eration 6: Am\'elioration du syst\`eme d'histoire / quete (branche)
\end{itemize}
\chapter{Bonus}
\begin{itemize}
	\item Ajout des achievements
	\item Beaucoup d'items
	\item Plus de classes
	\item Ladder solo
	\item Multijoueur
	\item Ladder multijoueur
	\item Animation enti\`ere du personnage
\end{itemize}

\end{document}
